% ============================================================
% v2 — Capitolo 5: Possibili evoluzioni
% Blocco III di IV. Budget: ~2 pagine.
% Le sette direzioni della v1 (§7.2), riorganizzate in tre gruppi
% per prossimità al lavoro svolto anziché in lista piatta.
% ============================================================
\chapter{Possibili evoluzioni}
\label{ch:evoluzioni}

Le direzioni che seguono sono ordinate per prossimità: prima quelle che
consolidano l'evidenza già raccolta, poi quelle che rimuovono i limiti
noti della piattaforma, infine quelle che estendono la domanda a un
terreno nuovo. Le prime due categorie sono realizzabili sul codice
esistente e su un budget di calcolo confrontabile con quello di questa
tesi, la terza no.

\section{Consolidare l'evidenza}
\label{sec:evo-evidenza}

\paragraph{Replicazione dei seed.}
È l'estensione più diretta e la più importante. Il design è stato
tagliato da tre seed appaiati per cella a uno per ragioni di calcolo
(\cref{sec:campagna}), e questa riduzione è la fonte principale di
cautela in tutta la tesi. Rieseguire la matrice $2\times2$ a $n{=}3$
trasformerebbe i casi di studio appaiati in stime di distribuzione e
metterebbe alla prova direttamente l'entità dell'interazione
architettura$\times$regime, che è oggi il risultato più forte e al tempo
stesso il meno replicato. Il costo è noto e lineare: otto run da tre
milioni di passi in aggiunta alle quattro esistenti.

\paragraph{Spiegare l'erosione.}
La degradazione entro la run controllata per livello
(\cref{tab:quartili}) è presente in tutte le celle e dimezzata sotto il
critic relazionale, e la tesi ne identifica due cause candidate senza
poterle separare. Due ablazioni mirate lo farebbero: congelare una
classe di agenti a metà addestramento isolerebbe la co-adattazione tra
classi, e spostare l'ancora dell'annealing di entropia isolerebbe lo
schedule di esplorazione. Entrambe sono modifiche di poche righe sulla
piattaforma esistente e girano al budget diagnostico del protocollo di
gate.

\paragraph{Selezione del checkpoint.}
Poiché i numeri finali sottostimano i picchi di metà addestramento in
ogni cella (\cref{sec:discussione}), un'analisi che valuti il
\emph{miglior} checkpoint anziché l'ultimo direbbe se l'ordinamento dei
task cambi anche il punto in cui la competenza culmina, e non solo
quello in cui atterra. I checkpoint intermedi sono già su disco: è
un'estensione della sola pipeline di valutazione.

\section{Rimuovere i limiti della piattaforma}
\label{sec:evo-piattaforma}

\paragraph{Generazione dei percorsi pedonali.}
Con più della metà dei percorsi pedonali sotto target su Town03
(\cref{sec:pedoni}), un planner consapevole della fattibilità ---
esplorazione più profonda delle ramificazioni, target adattivi per
regione di spawn, rinuncia esplicita quando la topologia locale non
sostiene la distanza richiesta --- è la correzione d'ambiente a maggior
leva: renderebbe interpretabili i tassi di successo medium e hard che
questa tesi è costretta a leggere solo entro livello, oggi vincolati da
un soffitto comune a tutte le celle.

\paragraph{Co-design di reward e curriculum.}
Il tuning della reward guidato da gate è girato per lo più su
configurazioni bloccate su easy (\cref{sec:minacce}), quindi lo shaping
del trunk può essere implicitamente sbilanciato verso i percorsi corti.
Ricalibrare i termini chiave sul mix completo quantificherebbe quanto
dell'effetto di regime sia mediato dalla calibrazione della reward, e
metterebbe alla prova la spiegazione avanzata per il fallimento di H1.

\section{Estendere la domanda}
\label{sec:evo-estensioni}

\paragraph{Estensioni architetturali sul critic.}
L'infrastruttura espone già encoder basati su attenzione (GAT) e la
normalizzazione del valore PopArt dietro flag di configurazione
(\cref{sec:critic}). Vista la forza dell'interazione osservata con
l'encoder del critic, queste non sono aggiunte esotiche ma le celle
successive naturali della stessa matrice: se sostituire un MLP con un
message passing a media sposta il \emph{luogo} dell'effetto di regime,
la domanda immediatamente successiva è se un'aggregazione pesata da
attenzione lo sposti ancora.

\paragraph{Curricula adattivi.}
Il gate di competenza di questa tesi è un sequenziatore fisso a due
soglie. Curricula che adattano le quote di livello in continuo --- per
esempio campionando in proporzione al progresso di apprendimento
misurato per livello --- metterebbero alla prova se l'effetto di
selezione dell'equilibrio sopravviva a un ordinamento a grana molto più
fine, o se dipenda proprio dalla discontinuità degli sblocchi.

\paragraph{Scalare la cornice.}
La copertura parallela dei task è stata fissata a $3{+}3$ agenti.
Scalare il mondo condiviso verso decine di agenti --- dove la struttura
relazionale del critic GNN dovrebbe contare di più --- trasforma la
cornice di coesistenza stessa in una variabile sperimentale, e sposta il
lavoro dalla caratterizzazione di un effetto alla sua legge di scala. È
l'unica direzione che richiede un'infrastruttura di calcolo
sostanzialmente diversa da quella usata qui
(\cref{app:riproducibilita}).
